\chapter{Bilangan Real}\label{ch:b1:reals}

Seluruh analisis bersandar pada satu sifat yang membedakan $\R$ dari
$\Q$: setiap \omterm{def:b1:logic:sets}{himpunan} tak kosong yang terbatas di atas mempunyai \omterm{def:b1:reals:bounds}{batas atas}
\emph{terkecil}. Bab ini menyatakannya secara tepat, menurunkan akibat
pertamanya --- \omterm{thm:b1:reals:archimedes}{sifat Archimedes}, \omterm{thm:b1:reals:floor}{fungsi lantai}, dan
kepadatan bilangan rasional maupun irasional --- lalu menyiapkan
kosakatanya (sup, inf, maks, min) yang dipakai terus-menerus mulai dari
\cref{ch:b1:seq}.

\section{Sifat batas atas}

\begin{definition}[Batas, sup dan inf]\label{def:b1:reals:bounds}
Misalkan $A \subseteq \R$ tak kosong. Bilangan real $M$ disebut \emph{batas
atas}\index{batas atas} $A$ bila $a \leq M$ untuk setiap $a \in A$;
lalu $A$ disebut \emph{terbatas di atas} bila ia mempunyai batas atas (serupa itu di
bawah, dengan batas bawah; \emph{terbatas} berarti keduanya). Adapun
\emph{maksimum} $A$ adalah batas atas yang termasuk $A$.

\emph{Supremum} $\sup A$\index{supremum} adalah batas atas terkecil
$A$, bila ada; sedangkan \emph{infimum} $\inf A$\index{infimum} adalah
batas bawah terbesarnya.
\end{definition}

\begin{theorem}[Aksioma kelengkapan $\R$]\label{thm:b1:reals:sup}
Di sini $\R$ adalah \omterm{def:b1:structures:field}{lapangan} terurut yang memuat $\Q$ dan yang di dalamnya \emph{setiap
\omterm{def:b1:logic:sets}{himpunan} bagian tak kosong yang terbatas di atas mempunyai \omterm{def:b1:reals:bounds}{supremum}}.\index{kelengkapan
R@kelengkapan $\R$}
\end{theorem}
\admitted

\begin{remark}
Kita mengambilnya sebagai aksioma pendefinisi $\R$; membangun sebuah modelnya (lewat
potongan Dedekind atau lewat barisan Cauchy bilangan rasional) beserta pembuktian
ketunggalannya memang jujur tetapi panjang, dan diserahkan untuk telaah lanjutan. Perhatikan
bahwa $\Q$ gagal memenuhi sifat itu: $\{x \in \Q : x^2 < 2\}$ terbatas
di atas tetapi tak mempunyai \omterm{def:b1:reals:bounds}{batas atas} terkecil \emph{di $\Q$} --- karena calonnya,
yaitu $\sqrt 2$, tidak ada di sana (\cref{ex:b1:logic:sqrt2}). Dengan berpindah ke
lawannya ($\sup(-A) = -\inf A$), setiap \omterm{def:b1:logic:sets}{himpunan} tak kosong yang terbatas di bawah
mempunyai \omterm{def:b1:reals:bounds}{infimum}.
\end{remark}

\begin{proposition}[Pencirian lewat $\varepsilon$]\label{prop:b1:reals:epsilon}
Misalkan $A \neq \emptyset$ terbatas di atas dan $s \in \R$. Maka $s =
\sup A$ jika dan hanya jika
\begin{enumerate}
  \item $s$ merupakan \omterm{def:b1:reals:bounds}{batas atas}: $\forall a \in A$, $a \leq s$; dan
  \item tak ada yang lebih kecil yang demikian: $\forall \varepsilon > 0$, $\exists a \in
        A$, $a > s - \varepsilon$.
\end{enumerate}
\end{proposition}

\begin{proof}
Jika $s = \sup A$: butir (1) berlaku menurut definisinya, dan untuk (2), $s -
\varepsilon < s$ bukan \omterm{def:b1:reals:bounds}{batas atas}, dan itu persis
keberadaan $a > s - \varepsilon$. Sebaliknya, (1) mengatakan $s$ sebuah
\omterm{def:b1:reals:bounds}{batas atas}; sedangkan (2) mengatakan tak ada $t < s$ yang menjadi \omterm{def:b1:reals:bounds}{batas atas} (ambil $\varepsilon
= s - t$): jadi $s$ yang terkecil.
\end{proof}

\begin{example}\label{ex:b1:reals:supexamples}
$\sup \intoo{0}{1} = 1$, yang tak tercapai (jadi tanpa maksimum);
$\sup \intcc{0}{1} = 1 = \max$. Untuk $A = \{1 - \frac 1n : n \in
\N^*\}$: $\sup A = 1$, yang tak tercapai; sedangkan $\inf A = \min A = 0$. Sebuah maksimum,
bila ada, sekaligus menjadi \omterm{def:b1:reals:bounds}{supremumnya}; dan seluruh maksud $\sup$ adalah menyediakan
pengganti ketika maksimumnya tidak ada.
\end{example}

\begin{omfigure}
\begin{tikzpicture}[scale=8.5]
  \draw[-Stealth] (-0.05,0) -- (1.3,0);
  \fill[omDef] (0,0) circle (0.35pt) node[below=2pt, font=\small]
    {$0$};
  \fill[omDef] (0.5,0) circle (0.35pt) node[below=2pt, font=\small]
    {$\tfrac12$};
  \fill[omDef] (0.6667,0) circle (0.35pt)
    node[below=2pt, font=\small] {$\tfrac23$};
  \fill[omDef] (0.75,0) circle (0.35pt)
    node[below=2pt, font=\small] {$\tfrac34$};
  \fill[omDef] (0.8,0) circle (0.35pt);
  \fill[omDef] (0.8333,0) circle (0.35pt);
  \fill[omDef] (0.8571,0) circle (0.35pt);
  \fill[omDef] (0.875,0) circle (0.35pt);
  \fill[omDef] (0.8889,0) circle (0.35pt);
  \fill[omDef] (0.9,0) circle (0.35pt);
  \draw[omThm, thick] (1,-0.015) -- (1,0.015);
  \node[above, omThm, font=\small] at (1,0.02) {$\sup A = 1$};
  \draw[omProp, very thick] (1,-0.04) -- (1.27,-0.04);
  \node[below, omProp, font=\small] at (1.13,-0.05)
    {batas atas $A$};
  \draw[Stealth-, gray] (0.93,0.03) -- (0.8,0.075)
    node[above left, font=\small, text=gray, inner sep=1pt]
    {ada $a > 1 - \varepsilon$};
\end{tikzpicture}

{\small \omterm{def:b1:logic:sets}{Himpunan} $A = \{1 - \frac1n : n \in \N^*\}$ pada garis
bilangan: titiknya menumpuk menuju $1$ tanpa pernah mencapainya. Setiap
bilangan $\geq 1$ merupakan \omterm{def:b1:reals:bounds}{batas atas} (yaitu sinar garisnya), dan tak ada yang lebih kecil
yang demikian, karena sebuah unsur $A$ masuk ke setiap \omterm{prop:b1:reals:intervals}{selang}
$\intoo{1 - \varepsilon}{1}$: jadi kedua klausa
\cref{prop:b1:reals:epsilon} dalam satu gambar. \omterm{def:b1:reals:bounds}{Supremumnya} adalah
titik ujung kiri sinar garis \omterm{def:b1:reals:bounds}{batas atas} itu --- dan aksioma kelengkapan
justru menjamin bahwa sinar garis itu selalu mempunyai titik ujung
kiri.}
\end{omfigure}

\begin{example}[Menghitung supremum dalam praktik]\label{ex:b1:reals:supcompute}
Dua latihan penuh atas \cref{prop:b1:reals:epsilon}.

\emph{\omterm{def:b1:logic:sets}{Himpunan} $A = \{x + \frac1x : x > 0\}$.} Untuk setiap $x > 0$,
$x + \frac1x - 2 = \frac{(\,\sqrt x - 1/\sqrt x\,)^2}{1} \geq 0$,
jadi $2$ merupakan batas bawah; dan $2 = 1 + \frac11 \in A$: sehingga
$\inf A = \min A = 2$, yang tercapai di $x = 1$. Di atas, $A$
tak terbatas ($x + \frac1x > x$ dapat melampaui sebarang $M$ menurut
\cref{thm:b1:reals:archimedes}): jadi $\sup A$ tidak ada di $\R$
(ia $+\infty$ di $\overline\R$).

\emph{\omterm{def:b1:logic:sets}{Himpunan} $B = \bigl\{\frac{m}{m + n} : m, n \in
\N^*\bigr\}$.} Setiap unsurnya terletak di $\intoo{0}{1}$, jadi $0$ dan
$1$ menjadi batasnya. Tak satu pun tercapai: karena $\frac{m}{m+n} = 1$ akan
memaksa $n = 0$. Untuk \omterm{def:b1:reals:bounds}{supremumnya}, bekukan $n = 1$ lalu biarkan $m$ tumbuh:
$\frac{m}{m+1} = 1 - \frac{1}{m+1} > 1 - \varepsilon$ begitu
$m + 1 > \frac1\varepsilon$ (menurut Archimedes): jadi $\sup B = 1$.
Secara simetris ($m = 1$, dengan $n$ besar), $\inf B = 0$. Inti gagasan
penutupnya: untuk memakukan sebuah \omterm{def:b1:reals:bounds}{supremum}, satu \emph{lintasan
berparameter tunggal} di dalam \omterm{def:b1:logic:sets}{himpunannya} sudah cukup --- di sini lintasan $n = 1$ --- dan
pencirian lewat $\varepsilon$ tak menuntut lebih dari itu.
\end{example}

\begin{example}[Cermin infimumnya]\label{ex:b1:reals:infmirror}
\omterm{def:b1:reals:bounds}{Infimum} mempunyai pencirian lewat $\varepsilon$ tersendiri, yang diperoleh
dari \cref{prop:b1:reals:epsilon} lewat $\inf A = -\sup(-A)$:
di sini $i = \inf A$ jika dan hanya jika $i$ membatasi $A$ dari bawah dan, untuk setiap
$\varepsilon > 0$, ada $a \in A$ dengan $a < i + \varepsilon$. Sebuah
latihan dengan kedua batasnya sekaligus: misalkan
\[
A = \Bigl\{(-1)^n + \frac1n : n \in \N^*\Bigr\}
= \Bigl\{0,\ \tfrac32,\ -\tfrac23,\ \tfrac54,\ -\tfrac45,\
\dots\Bigr\} .
\]
Indeks genap memberikan $1 + \frac1n \leq \frac32$, dengan kesamaan di
$n = 2$: dan karena nilai berindeks ganjil pun $\leq 0 <
\frac32$, kita peroleh $\sup A = \max A = \frac32$. Sedangkan indeks ganjil memberikan
$-1 + \frac1n > -1$, yang menurun menuju $-1$: jadi setiap unsur
$A$ bernilai $> -1$, dan $-1 + \varepsilon$ dikalahkan oleh $-1 +
\frac1n$ untuk $n > \frac1\varepsilon$ yang ganjil: sehingga $\inf A = -1$, yang tak
tercapai. Satu \omterm{def:b1:logic:sets}{himpunan}, dengan keempat perilakunya terpampang: sebuah \omterm{def:b1:reals:bounds}{supremum}
yang menjadi maksimum, dan sebuah \omterm{def:b1:reals:bounds}{infimum} yang bukan minimum.
\end{example}

\begin{remark}[Jebakan yang lazim dengan sup dan inf]
Empat kesalahan menyumbang sebagian besar nilai yang hilang. (i) \emph{Mengacaukan
$\sup$ dan $\max$}: karena $\sup A$ tak harus termasuk $A$; jadi tulis $\max$
hanya setelah sebuah unsur $A$ yang menjadi \omterm{def:b1:reals:bounds}{batas atas} ditunjukkan.
(ii) \emph{Meneruskan ketaksamaan tegas ke \omterm{def:b1:reals:bounds}{supremumnya}}: jika $a <
b$ untuk setiap $a \in A$, kita hanya boleh menyimpulkan $\sup A \leq b$ ---
dengan saksi $A = \intoo{0}{1}$, $b = 1$. (iii) \emph{Menulis $\sup A$
sebelum memeriksa kesahihannya}: karena lambang itu menuntut $A$ tak kosong dan
terbatas di atas (\cref{met:b1:reals:supproofs}); jadi $\sup \emptyset$
dan $\sup \N$ tak terdefinisi di $\R$ (adapun kesepakatan pada
$\overline\R$ adalah tindakan tersendiri yang eksplisit). (iv) \emph{Operasi
\omterm{def:b1:logic:sets}{himpunan}}: di sini $\sup(A \cup B) = \max(\sup A, \sup B)$ selalu berlaku, tetapi
tak ada yang berlaku umum untuk $A \cap B$ --- karena ia mungkin kosong, dan
bahkan ketika tidak, $\sup(A \cap B)$ dapat jauh di bawah
$\min(\sup A, \sup B)$: ambil $A = \{0, 2\}$ dan $B = \{0, 3\}$,
yang di situ $\sup(A \cap B) = 0$.
\end{remark}

\begin{example}[Himpunan hingga mempunyai maksimum --- lema yang dipakai diam-diam]\label{ex:b1:reals:finitemax}
Setiap $F \subseteq \R$ yang hingga dan tak kosong mempunyai maksimum (dan
minimum). Induksi pada banyaknya unsurnya: singleton
$\{a\}$ mempunyai $\max = a$; lalu jika klaimnya berlaku untuk \omterm{def:b1:logic:sets}{himpunan} beranggota $n$
dan $F$ beranggota $n + 1$ unsur, pilih sebarang $a \in F$: \omterm{def:b1:logic:sets}{himpunan} $F
\setminus \{a\}$ mempunyai maksimum $m$, dan $\max F$ adalah $m$ bila $a
\leq m$, dan $a$ bila tidak. Tak ada kelengkapan yang terlibat --- ini murni
urutan ditambah induksi, yang sudah sah di $\Q$ --- namun lemanya
pantas mendapat satu \omterm{def:b1:logic:statement}{pernyataan} yang jujur karena bukti berikutnya memanggilnya
diam-diam: baik konstruksi \omterm{thm:b1:reals:floor}{lantai} di bawah (``\omterm{def:b1:logic:sets}{himpunan} bilangan bulat
yang terperangkap pada rentang hingga mempunyai unsur terbesar''), setiap
batas $\max(\abs{u_0}, \dots, \abs{u_{N-1}}, \dots)$ pada
\cref{ch:b1:seq}, maupun setiap ``ambil yang terbesar di antara $\delta$ yang berhingga banyak''
pada \cref{ch:b1:continuity}. Adapun \omterm{def:b1:logic:sets}{himpunan} tak hingga adalah tempat
maksimum mati dan \omterm{def:b1:reals:bounds}{supremum} mengambil alih: bab ini ada demi kasus yang
tak hingga itu.
\end{example}

\begin{theorem}[Sifat Archimedes]\label{thm:b1:reals:archimedes}
Untuk setiap $x \in \R$ ada $n \in \N$ dengan $n > x$. Setara dengan itu:
untuk setiap $\varepsilon > 0$ dan $y > 0$, ada kelipatan $n\varepsilon$
yang melampaui $y$.\index{sifat Archimedes}
\end{theorem}

\begin{proof}
Andaikan tidak: maka ada $x$ yang menjadi \omterm{def:b1:reals:bounds}{batas atas} $\N$. Lalu $s = \sup \N$
ada (\cref{thm:b1:reals:sup}). Menurut
\cref{prop:b1:reals:epsilon}~(2) dengan $\varepsilon = 1$, ada $n
\in \N$ dengan $n > s - 1$; tetapi lalu $n + 1 \in \N$ dan $n + 1 > s$,
yang bertentangan dengan $s$ sebagai \omterm{def:b1:reals:bounds}{batas atas}. Untuk bentuk keduanya,
misalkan $\varepsilon > 0$ dan $y > 0$: bentuk pertamanya yang diterapkan pada $x
= \frac{y}{\varepsilon}$ menghasilkan $n \in \N$ dengan $n >
\frac{y}{\varepsilon}$, lalu mengalikannya dengan $\varepsilon > 0$
(yang mengawetkan ketaksamaan tegas) memberikan $n\varepsilon > y$.
Sebaliknya, bentuk keduanya dengan $\varepsilon = 1$ dan $y = x$
memulihkan yang pertama untuk $x > 0$, dan $n = 1$ menangani $x \leq 0$:
jadi kedua \omterm{def:b1:logic:statement}{pernyataannya} setara secara tegas.
\end{proof}

\begin{example}[Archimedes dalam praktik]\label{ex:b1:reals:archimedeswork}
Tiga pemakaian langsung, yang terus-menerus diperlukan nanti. (i) \emph{Tak ada
bilangan real positif yang berada di bawah setiap $\frac1n$}: karena bila $0 < \varepsilon$,
pilihlah $n > \frac1\varepsilon$; maka $\frac1n < \varepsilon$. Dengan
kata lain, $\R$ tak memuat infinitesimal --- dan ungkapan tak resmi
``$\frac1n$ menjadi sekecil-kecilnya'' persis merupakan teorema ini.
(ii) \emph{Ambang yang eksplisit}: seberapa besar $n$ harus dibuat agar
$\frac{1}{n^2} \leq 10^{-6}$? Cukuplah $n \geq 10^3$
--- Archimedes menjamin bahwa $n$ semacam itu ada, dan aljabarnya
menemukan letaknya. (iii) \emph{Pangkat mengalahkan sebarang batas}: karena $2^n \geq n +
1$ (lewat induksi), maka untuk setiap $M$ ada pangkat $2$ yang melampaui $M$:
yaitu pertumbuhan geometri yang dipakai bagi bilangan diadik pada
\cref{exo:b1:reals:8}. Inti gagasan penutupnya: \omterm{thm:b1:reals:archimedes}{sifat Archimedes}
adalah izin di balik setiap frasa berbentuk ``ambil
$n$ yang cukup besar'' --- jadi mulai sekarang kita memakai frasa itu dengan bebas,
dan contoh ini adalah pembenarannya sekali untuk selamanya.
\end{example}

\begin{theorem}[Fungsi lantai]\label{thm:b1:reals:floor}
Untuk setiap $x \in \R$ ada tepat satu bilangan bulat, yaitu
\emph{lantai}\index{fungsi lantai}nya $\lfloor x \rfloor$, dengan
\[
\lfloor x \rfloor \leq x < \lfloor x \rfloor + 1 .
\]
\end{theorem}

\begin{proof}
\emph{Keberadaannya.} \omterm{def:b1:logic:sets}{Himpunan} $E = \{k \in \Z : k \leq x\}$ tak kosong:
karena menurut \cref{thm:b1:reals:archimedes} ada $m \in \N$ dengan $m > -x$,
sehingga $-m < x$, jadi $-m \in E$. Ia terbatas di atas (oleh sebarang bilangan bulat
$n > x$, yang ada karena alasan yang sama), sehingga, sebagai \omterm{def:b1:logic:sets}{himpunan} bilangan
bulat yang terperangkap pada rentang hingga $\intint{-m}{n}$, ia mempunyai
unsur terbesar $k = \max E$. Maka $k \leq x$, dan $k + 1 \notin E$
berarti $x < k + 1$.

\emph{Ketunggalannya.} Jika $k$ dan $k'$ sama-sama memenuhi ketaksamaannya,
maka $k \leq x < k' + 1$ memberikan $k \leq k'$, dan secara simetris $k'
\leq k$.
\end{proof}

\begin{example}[Lantai dalam praktik]\label{ex:b1:reals:floorpractice}
$\lfloor 3.7 \rfloor = 3$, $\lfloor 5 \rfloor = 5$, dan $\lfloor
-3.7 \rfloor = -4$: jadi \omterm{thm:b1:reals:floor}{lantainya} turun \emph{ke bawah}, bukan menuju $0$.
Dua akibat ketunggalan pada \cref{thm:b1:reals:floor} yang
akan kita pakai diam-diam. Pertama, untuk $n \in \Z$,
\[
\lfloor x + n \rfloor = \lfloor x \rfloor + n ,
\]
karena $\lfloor x \rfloor + n$ adalah bilangan bulat yang memenuhi kedua
ketaksamaan pendefinisi bagi $x + n$ --- dan hanya satu bilangan bulat yang demikian.
Kedua, $\lfloor \, \cdot \, \rfloor$ tidak turun: karena jika $x \leq
y$ maka $\lfloor x \rfloor \leq x \leq y < \lfloor y \rfloor + 1$,
dan bilangan bulat yang $< \lfloor y \rfloor + 1$ pastilah $\leq \lfloor y
\rfloor$. Awas, bagaimanapun, bahwa $\lfloor 2x \rfloor \neq 2\lfloor
x \rfloor$ pada umumnya: karena $x = 0.6$ memberikan $\lfloor 1.2 \rfloor = 1
\neq 0 = 2\lfloor 0.6 \rfloor$.

Yang \emph{memang} benar adalah sebuah kesamaan terkerjakan yang layak disimpan (kesamaan
Hermite, dalam kasusnya yang paling sederhana): untuk setiap $x$ real,
\[
\lfloor x \rfloor + \Bigl\lfloor x + \frac12 \Bigr\rfloor
= \lfloor 2x \rfloor .
\]
Tulis $x = \lfloor x\rfloor + u$ dengan $u \in \intco{0}{1}$ lalu
pisahkan dua kasus. Jika $u < \frac12$: ruas kirinya menjadi $\lfloor
x\rfloor + \lfloor x\rfloor = 2\lfloor x\rfloor$, dan $2x =
2\lfloor x\rfloor + 2u$ dengan $2u \in \intco{0}{1}$, sehingga ruas
kanannya juga $2\lfloor x\rfloor$. Jika $u \geq \frac12$: ruas
kirinya menjadi $\lfloor x\rfloor + (\lfloor x\rfloor + 1)$, dan $2u \in
\intco{1}{2}$ membuat ruas kanannya $2\lfloor x\rfloor + 1$. Inti
gagasan penutupnya: $\lfloor x + \frac12\rfloor$ adalah
\emph{pembulatan} $x$ ke bilangan bulat terdekat, sehingga kesamaannya
mengatakan bahwa \omterm{thm:b1:reals:floor}{lantai} ditambah pembulatan sama dengan \omterm{thm:b1:reals:floor}{lantai} dari lipat duanya --- dan
pemilahan kasus atas bagian pecahannya $u$ adalah teknik
baku di balik setiap kesamaan \omterm{thm:b1:reals:floor}{lantai}
(\cref{exo:b1:reals:2,exo:b1:reals:3} berjalan di atasnya juga).
\end{example}

\begin{theorem}[Kepadatan $\Q$ dan $\R \setminus \Q$]\label{thm:b1:reals:density}
Di antara sebarang dua bilangan real $x < y$ terdapat sebuah bilangan rasional dan sebuah
bilangan irasional.\index{kepadatan}
\end{theorem}

\begin{proof}
\emph{Sebuah bilangan rasional.} Menurut \cref{thm:b1:reals:archimedes}, pilih $n \in
\N^*$ dengan $n > \frac{1}{y - x}$, sehingga $ny - nx > 1$. Misalkan $m = \lfloor
nx \rfloor + 1$. Di satu sisi, $nx < \lfloor nx \rfloor + 1 = m$
(\cref{thm:b1:reals:floor}); di sisi lain, $m = \lfloor nx \rfloor +
1 \leq nx + 1 < ny$. Dengan \omterm{def:b1:arith:divides}{membaginya} dengan $n$: $x < \frac mn < y$.

\emph{Sebuah bilangan irasional.} Terapkan butir sebelumnya pada pasangan $x -
\sqrt 2 < y - \sqrt 2$: maka ada bilangan rasional $q$ di antaranya, dan lalu
$q + \sqrt 2 \in \intoo{x}{y}$ bersifat irasional (karena seandainya $q + \sqrt 2$
rasional, maka $\sqrt 2$ pun rasional).
\end{proof}

\begin{example}[Menjalankan bukti kepadatannya]\label{ex:b1:reals:densityrun}
Buktinya adalah sebuah algoritma; mari kita jalankan pada $x = 1.414$ dan
$y = \sqrt 2$. Karena $1.4142^2 = 1.99996164 < 2$, kita punya $\sqrt 2
> 1.4142$, sehingga $y - x > 0.0002$ dan $\frac{1}{y - x} < 5000$: jadi
pilihan $n = 5000$ sah. Lalu $nx = 7070$, sehingga $m = \lfloor
7070 \rfloor + 1 = 7071$, dan bilangan rasional yang dihasilkannya adalah
\[
\frac{m}{n} = \frac{7071}{5000} = 1.4142,
\qquad
1.414 < 1.4142 < \sqrt 2 .
\]
Inti gagasan penutupnya: buktinya hanya memerlukan $n$ yang sedikit lebih besar daripada
$\frac{1}{y-x}$, dan ia mengembalikan kelipatan $\frac 1n$ yang pertama
di luar $x$. Jadi kepadatan bukanlah mukjizat yang abstrak --- ia
pembagian panjang yang menyamar, yaitu tema yang dikembangkan panjang lebar pada soal akhir
pekan (\cref{pb:b1:reals:1}).
\end{example}

\begin{remark}[Di mana kelengkapan dipakai berikutnya]
\Cref{thm:b1:reals:sup} adalah satu-satunya aksioma tak aljabar buku
ini, dan setiap teorema keberadaan dalam analisis hanyalah aksioma itu
yang berganti pakaian: teorema kekonvergenan monoton
(\cref{ch:b1:seq}), teorema Bolzano--Weierstrass
(\cref{ch:b1:topology}), teorema nilai antara dan teorema nilai
ekstrem (\cref{ch:b1:continuity}), serta definisi
integral itu sendiri sebagai \omterm{def:b1:reals:bounds}{supremum} jumlah bawah (\cref{ch:b1:integration}).
Jilid Tahun ke-3 membangun teori ukuran dan ruang Hilbert di atas
aksioma tunggal yang sama. Jadi ketika sebuah bukti pada bab berikutnya menghasilkan
bilangan real entah dari mana, carilah \omterm{def:b1:reals:bounds}{supremum} yang tersembunyi.
\end{remark}

\begin{remark}[Antara kediskretan dan kepadatan]
\omterm{def:b1:logic:sets}{Himpunan} $\Z$ dan $\Q$ duduk pada dua ujung yang berlawanan di dalam $\R$: di sekeliling
setiap bilangan bulat ada celah sepanjang $1$ yang tak memuat bilangan bulat lain
(kediskretan --- dan itulah yang membuat \omterm{thm:b1:reals:floor}{lantai} terdefinisi dengan baik),
sedangkan di antara dua bilangan real mana pun ada tak hingga banyak bilangan rasional
(kepadatan). Menariknya, untuk \emph{\omterm{def:b1:structures:subgroup}{subgrup} aditif} $\R$
tak ada apa pun di antaranya: \cref{exo:b1:reals:9} membuktikan bahwa
\omterm{def:b1:structures:subgroup}{subgrup} semacam itu entah berbentuk $\alpha\Z$ (diskret) atau
padat --- yaitu dikotomi yang menggerakkan kepadatan $\{\sin n\}$ pada
\cref{ch:b1:seq} dan monster konstruktif pada
\cref{pb:b1:continuity:1}. Adapun \omterm{def:b1:logic:sets}{himpunan} yang umum, tentu saja, mencampur kedua
perilakunya dengan bebas: $\Z \cup \Q\cap\intcc{0}{1}$ diskret jauh
di sana dan padat di tengahnya.
\end{remark}

\begin{method}[Membuktikan kesamaan dengan sup dan inf]\label{met:b1:reals:supproofs}
Untuk membuktikan $\sup A = s$: periksa bahwa $s$ membatasi $A$ dari atas, lalu hasilkan,
untuk setiap $\varepsilon > 0$ (atau untuk barisan $\varepsilon = \frac
1n$), sebuah unsur $A$ di atas $s - \varepsilon$. Untuk membandingkan \omterm{def:b1:reals:bounds}{supremum},
pakailah: $A \subseteq B \implies \sup A \leq \sup B$; dan untuk setiap $a, b$:
$\sup(A + B) = \sup A + \sup B$, dengan $A + B = \{a + b\}$
(\cref{exo:b1:reals:5}). Jangan pernah menulis $\sup A$ sebelum mengetahui bahwa $A$
tak kosong dan terbatas di atas.
\end{method}

\section{Selang}

\begin{proposition}[Pencirian selang]\label{prop:b1:reals:intervals}
\omterm{def:b1:logic:sets}{Himpunan} bagian $I \subseteq \R$ merupakan \emph{selang}\index{selang} (yaitu salah
satu jenis yang sudah dikenal $\intoo{a}{b}$, $\intcc{a}{b}$, $\intco{a}{b}$,
$\intoc{a}{b}$, setengah garis, $\R$, $\emptyset$, singleton) jika dan
hanya jika ia \emph{konveks}:
\[
\forall x, y \in I,\ \forall z \in \R, \quad
x \leq z \leq y \implies z \in I .
\]
\end{proposition}

\begin{proof}
Setiap jenis yang didaftar jelas konveks. Sebaliknya, misalkan $I$ konveks dan
tak kosong. Tulis $a = \inf I$ bila $I$ terbatas di bawah, dan $a =
-\infty$ bila tidak; demikian pula $b = \sup I$ atau $+\infty$. Kita klaim $\intoo{a}{b}
\subseteq I \subseteq \intcc{a}{b}$ (dengan kesepakatan yang kasatmata di
$\pm\infty$). Inklusi keduanya adalah definisi batasnya. Untuk
yang pertama, misalkan $z \in \intoo{a}{b}$: karena $z > a$, maka $z$ bukan
batas bawah (atau $a = -\infty$), sehingga ada $x \in I$ dengan $x < z$;
demikian pula ada $y \in I$ dengan $y > z$; lalu kekonveksannya menempatkan $z \in I$.

Tinggal membaca jenisnya dari inklusi ganda
$\intoo{a}{b} \subseteq I \subseteq \intcc{a}{b}$: \omterm{def:b1:logic:sets}{himpunan} yang
terapit di antara \omterm{prop:b1:reals:intervals}{selang} terbuka dan penutupnya hanya berbeda dari
$\intoo{a}{b}$ oleh hadir atau tidaknya titik ujung yang
hingga. Secara eksplisit: jika $a, b \in \R$, keempat kemungkinan
bagi $(a \in I,\ b \in I)$ memberikan $\intoo{a}{b}$, $\intco{a}{b}$,
$\intoc{a}{b}$, $\intcc{a}{b}$ (termasuk kasus merosot
$a = b$: yaitu singleton bila $a \in I$); jika $a = -\infty$ dan $b \in
\R$, kita peroleh $\intoo{-\infty}{b}$ atau $\intoc{-\infty}{b}$;
simetris dengan itu untuk $a \in \R$, $b = +\infty$; dan $a = -\infty$,
$b = +\infty$ memberikan $I = \R$. Setiap kasusnya ada di daftar: selesai.
\end{proof}

\begin{remark}[Mengapa kekonveksan adalah uji yang tepat]
Proposisi itu mengubah definisi yang \emph{geometris} (yaitu daftar
sepuluh bentuk) menjadi uji \emph{logika satu baris}, dan ujinyalah
yang benar-benar kita pakai: untuk membuktikan bahwa sebuah \omterm{def:b1:logic:sets}{himpunan} merupakan \omterm{prop:b1:reals:intervals}{selang},
jangan pernah mengejar bentuk mana di antara kesepuluhnya --- periksalah kekonveksannya
lalu biarkan proposisinya yang memilah jenisnya. Teorema nilai
antara pada \cref{ch:b1:continuity} akan dinyatakan persis
dengan cara ini (``peta kontinu sebuah \omterm{prop:b1:reals:intervals}{selang} adalah
\omterm{prop:b1:reals:intervals}{selang}''), dan buktinya menghasilkan kekonveksannya, bukan
bentuknya.
\end{remark}

\begin{remark}[Garis real yang diperluas]
Nyaman sekali menambahkan dua lambang lalu bekerja di $\overline\R = \R
\cup \{-\infty, +\infty\}$, dengan kesepakatan $\sup A = +\infty$
bila $A$ tak terbatas di atas dan $\sup \emptyset = -\infty$. Maka
\emph{setiap} \omterm{def:b1:logic:sets}{himpunan} bagian $\R$ mempunyai \omterm{def:b1:reals:bounds}{supremum} di $\overline\R$ --- yaitu
kenyamanan notasi yang dipakai dengan bebas untuk limit pada \cref{ch:b1:seq}.
\end{remark}

\begin{example}[Berhitung di $\overline\R$]\label{ex:b1:reals:extendedcompute}
Dengan kesepakatan itu berlaku: $\sup \Z = +\infty$, $\inf \Z =
-\infty$; untuk $A = \{n + (-1)^n n : n \in \N\} = \{0, 4, 0, 8,
\dots\} \cup \{0\}$, $\sup A = +\infty$ (karena suku genapnya $2n$
tak terbatas) dan $\inf A = \min A = 0$; dan $\sup\emptyset =
-\infty \leq \inf\emptyset = +\infty$ --- yaitu satu-satunya \omterm{def:b1:logic:sets}{himpunan} yang
\omterm{def:b1:reals:bounds}{supremumnya} \emph{lebih kecil} daripada \omterm{def:b1:reals:bounds}{infimumnya}, sebuah pengingat bahwa
kesepakatannya dipilih agar $\sup$ naik dan $\inf$
turun terhadap inklusi:
\[
A \subseteq B \implies \sup A \leq \sup B
\quad\text{dan}\quad \inf A \geq \inf B ,
\]
yang kini sah tanpa peringatan ketakkosongan. Yang \emph{tidak}
disediakan kesepakatan itu adalah aritmetikanya: $+\infty + (-\infty)$ dan $0
\times (+\infty)$ tetap tak terdefinisi, dan setiap manipulasi aljabar
atas \omterm{def:b1:reals:bounds}{supremum} harus lebih dulu memeriksa bahwa ia tak pernah membentuk keduanya.
Jadi garis yang diperluas itu pembukuan, bukan sistem bilangan.
\end{example}

\begin{example}[Supremum yang lolos dari $\Q$]\label{ex:b1:reals:supescape}
Kembalilah ke \omterm{def:b1:logic:sets}{himpunan} pada catatan pembuka, $A = \{x \in \Q : x^2 <
2\}$, lalu hitung \omterm{def:b1:reals:bounds}{supremumnya} \emph{di $\R$}. Ia tak kosong
(karena $1 \in A$) dan terbatas di atas oleh $1.5$ (karena jika $x > 1.5$ maka $x^2 >
2.25 > 2$), sehingga $s = \sup A$ ada. Kita klaim $s = \sqrt 2$ (yaitu
bilangan real yang dibangun pada \cref{exo:b1:reals:12}). Batas atasnya: setiap
$a \in A$ memenuhi $a < \sqrt2$ --- untuk $a \leq 0$ ini
jelas, dan untuk $a > 0$, $a \geq \sqrt2$ akan memberikan $a^2 \geq 2$.
Tak ada yang lebih kecil yang berhasil: diberikan $t < \sqrt2$, kepadatan
(\cref{thm:b1:reals:density}) memasok bilangan rasional $q$ dengan
$\max(1, t) < q < \sqrt 2$, dan lalu $q^2 < 2$, sehingga $q \in A$
melampaui $t$. Menurut \cref{prop:b1:reals:epsilon}, $s = \sqrt2
\notin \Q$. Inti gagasan penutupnya: \omterm{def:b1:reals:bounds}{supremum} sebuah \omterm{def:b1:logic:sets}{himpunan}
bilangan rasional tak harus rasional --- dan kelengkapan justru
janji bahwa $\R$, tidak seperti $\Q$, tak pernah membiarkan sebuah \omterm{def:b1:reals:bounds}{supremum} lolos;
jadi contoh ini adalah catatan pembuka bab ini, yang kini terbukti
alih-alih sekadar ditunjuk.
\end{example}

\begin{remark}[Cakrawala di dalam jilid ini]
Ketiga perkakas bab ini mempunyai karier yang berbeda-beda di depan. Adapun
\omterm{def:b1:reals:bounds}{supremum} menjalankan paruh analisisnya: limit monoton
(\cref{ch:b1:seq}), definisi integral itu sendiri
(\cref{ch:b1:integration}), dan, pada geometri
\cref{ch:b1:euclid}, jarak dari sebuah titik ke sebuah subruang ---
yaitu \omterm{def:b1:reals:bounds}{infimum} yang diubah proyeksi ortogonal menjadi minimum. Adapun
\omterm{thm:b1:reals:floor}{fungsi lantai} kembali di mana pun yang diskret bertemu yang
kontinu: ekspansi angka (soal akhir pekan bab ini),
hampiran sarang merpati Dirichlet
(\cref{pb:b1:derivative:1}), perbandingan integral atas jumlah
(\cref{ch:b1:series}). Sedangkan argumen kepadatan naik pangkat menjadi sebuah metode
pada \cref{ch:b1:continuity}: kesamaan antara fungsi kontinu
hanya perlu diperiksa pada $\Q$ --- dan separuh persamaan fungsional
Cauchy (\cref{pb:b1:continuity:1}) persis merupakan langkah itu. Jadi ketika
ragu tentang dari mana sebuah bukti dalam jilid ini memperoleh \omterm{def:b1:logic:statement}{pernyataan}
keberadaannya, jawabannya hampir selalu: bab ini.
\end{remark}

\section{Latihan}

\begin{exercise}[$\star$]\label{exo:b1:reals:1}
Tentukan (beserta buktinya) sup, inf, maks, min --- bila ada --- dari:
\[
A = \Bigl\{\frac{1}{n} : n \in \N^*\Bigr\},
\qquad
B = \Bigl\{\frac{(-1)^n n}{n+1} : n \in \N\Bigr\},
\qquad
C = \{x \in \R : x^2 < 3\}.
\]
\end{exercise}

\begin{exercise}[$\star$]\label{exo:b1:reals:2}
Buktikan bahwa untuk setiap $x, y \in \R$: $\lfloor x \rfloor + \lfloor y
\rfloor \leq \lfloor x + y \rfloor \leq \lfloor x \rfloor + \lfloor y
\rfloor + 1$, dan bahwa kedua batasnya tercapai.
\end{exercise}

\begin{exercise}[$\star$]\label{exo:b1:reals:3}
Buktikan bahwa untuk setiap $x \in \R$ dan $n \in \N^*$:
$\Bigl\lfloor \frac{\lfloor nx \rfloor}{n} \Bigr\rfloor = \lfloor x
\rfloor$.
\end{exercise}

\begin{exercise}[$\star$]\label{exo:b1:reals:4}
Misalkan $A \subseteq B$ \omterm{def:b1:logic:sets}{himpunan} bagian tak kosong $\R$, dengan $B$ terbatas. Buktikan
$\inf B \leq \inf A \leq \sup A \leq \sup B$.
\end{exercise}

\begin{exercise}[$\star\star$]\label{exo:b1:reals:5}
Untuk $A, B \subseteq \R$ yang tak kosong dan terbatas, definisikan $A + B = \{a + b : a
\in A,\ b \in B\}$ dan $-A = \{-a : a \in A\}$. Buktikan:
\[
\sup(A + B) = \sup A + \sup B,
\qquad
\sup(-A) = -\inf A .
\]
\end{exercise}

\begin{exercise}[$\star\star$]\label{exo:b1:reals:6}
Misalkan $f, g \colon E \to \R$ fungsi yang terbatas. Buktikan
\[
\sup_{x \in E}\, \bigl(f(x) + g(x)\bigr) \leq \sup_{x \in E} f(x) +
\sup_{x \in E} g(x),
\]
lalu berikan contoh yang ketaksamaannya tegas. Mengapa hal ini tidak
bertentangan dengan \cref{exo:b1:reals:5}?
\end{exercise}

\begin{exercise}[$\star\star$]\label{exo:b1:reals:7}
Buktikan bahwa $\sqrt 2 + \sqrt 3$ irasional. \emph{(Kuadratkan lalu
pakai keirasionalan $\sqrt 6$, yang akan dibuktikan lewat
\cref{exo:b1:arith:7}.)}
\end{exercise}

\begin{exercise}[$\star\star$]\label{exo:b1:reals:8}
Buktikan bahwa \omterm{def:b1:logic:sets}{himpunan} $D = \bigl\{\frac{m}{2^n} : m \in \Z,\ n \in
\N\bigr\}$ berisi bilangan rasional diadik bersifat padat di $\R$: yaitu di antara dua
bilangan real mana pun terdapat sebuah bilangan rasional diadik.
\end{exercise}

\begin{exercise}[$\star\star\star$]\label{exo:b1:reals:9}
Misalkan $G$ sebuah \omterm{def:b1:structures:subgroup}{subgrup} $(\R, +)$ dengan $G \neq \{0\}$. Tulis $\alpha =
\inf\,(G \cap \intoo{0}{+\infty})$. Buktikan:
\begin{enumerate}
  \item jika $\alpha > 0$, maka $G = \alpha\Z$;
  \item jika $\alpha = 0$, maka $G$ padat di $\R$.
\end{enumerate}
Lalu simpulkan bahwa $\Z + \sqrt 2\,\Z$ padat di $\R$.
\end{exercise}

\begin{exercise}[$\star\star\star$]\label{exo:b1:reals:10}
Untuk $A, B$ \omterm{def:b1:logic:sets}{himpunan} tak kosong berisi bilangan real positif, tulis $AB = \{ab : a \in
A, b \in B\}$. Buktikan $\sup(AB) = \sup A \cdot \sup B$ (pada kasus terbatas),
lalu tunjukkan lewat contoh bahwa kepositifannya penting.
\end{exercise}

\begin{exercise}[$\star\star$]\label{exo:b1:reals:11}
Untuk $A \subseteq \R$ yang tak kosong dan terbatas, definisikan
\emph{diameternya}
\[
\operatorname{diam} A = \sup\,\{\abs{a - a'} : a, a' \in A\} .
\]
Buktikan bahwa $\operatorname{diam} A = \sup A - \inf A$, dan bahwa
$\intcc{\inf A}{\sup A}$ adalah \omterm{prop:b1:reals:intervals}{selang} tertutup terkecil
yang memuat $A$.
\end{exercise}

\begin{exercise}[$\star\star\star$]\label{exo:b1:reals:12}
Misalkan $y > 0$ dan $E = \{x \geq 0 : x^2 \leq y\}$. Buktikan bahwa $E$
tak kosong dan terbatas di atas, dan bahwa $s = \sup E$ memenuhi $s^2 =
y$ \emph{(singkirkan $s^2 < y$ dan $s^2 > y$ dengan menunjukkan, pada masing-masing
kasus, sebuah $h > 0$ yang kecil yang bertentangan dengan definisi
\omterm{def:b1:reals:bounds}{supremumnya})}. Lalu simpulkan bahwa setiap $y > 0$ mempunyai akar kuadrat tunggal
$\sqrt y > 0$ dan bahwa $y \mapsto \sqrt y$ naik pada
$\intoo{0}{+\infty}$.
\end{exercise}

\section{Soal: Ekspansi angka dan irama bilangan rasional}

\begin{problem}[{Soal akhir pekan --- ekspansi $b$-adik:
keberadaan, ketunggalan, dan keperiodikan yang mencirikan $\Q$}]
\label{pb:b1:reals:1}
Setiap bilangan real di $\intco{0}{1}$ mempunyai ekspansi angka dalam setiap basis
$b \geq 2$; ekspansinya tunggal begitu untaian angka $b - 1$ di
belakang dilarang; dan ia periodik pada akhirnya tepat
ketika bilangannya rasional. Soal ini membuktikan ketiga fakta itu
dari aksioma kelengkapan saja --- tanpa barisan, tanpa deret:
hanya dengan \omterm{def:b1:reals:bounds}{supremum}, \omterm{thm:b1:reals:archimedes}{sifat Archimedes} dan \omterm{thm:b1:reals:floor}{fungsi lantai}
--- lalu menutup dengan argumen diagonal Cantor dalam wujud angka.
Di sepanjang soal ini, $b \geq 2$ adalah bilangan bulat tetap (yaitu \emph{basisnya}), sebuah
\emph{angka} adalah unsur $\intint{0}{b-1}$, dan sebuah untaian
angka $(d_n)_{n \geq 1}$ disebut \emph{sejati} bila ia tidak
akhirnya sama dengan $b - 1$ (yakni: untuk setiap $N$ ada $n >
N$ dengan $d_n \leq b - 2$).\index{ekspansi b-adik@ekspansi $b$-adik}\index{ekspansi desimal}

\medskip
\noindent\textbf{Bagian I --- Angka dengan tangan.} Pembagian panjang $p$
oleh $q$ dalam basis $b$: kalikan sisa berjalannya dengan $b$, bagi
dengan $q$, catat hasil baginya sebagai angka berikutnya, lalu simpan sisanya.
\begin{enumerate}
  \item Dalam basis $10$, jalankan algoritmanya pada $\frac 18$ dan pada
        $\frac 17$, sambil mencatat pada setiap langkahnya angkanya \emph{beserta}
        sisanya. Periksa bahwa sisa untuk $\frac 17$
        berdaur melewati $1, 3, 2, 6, 4, 5$ dan bahwa angka
        $142857$ lalu berulang selamanya.
  \item Hitung ekspansi basis $2$ dari $\frac 13$ dan dari
        $\frac{5}{16}$, dan ekspansi basis $3$ dari $\frac 12$.
        Amati: satu bilangan berakhir, dua lainnya berulang ---
        dan $\frac 12$, yang begitu jinak dalam basis $10$, berulang selamanya dalam
        basis $3$.
  \item Untuk $x = \frac pq \in \intco{0}{1}$ dalam bentuk paling sederhana, tunjukkan
        bahwa angka yang dihasilkan algoritmanya akhirnya semuanya
        $0$ jika dan hanya jika sisa $b^N p \bmod q$
        lenyap untuk suatu $N$, jika dan hanya jika $q$ \omterm{def:b1:arith:divides}{membagi} suatu
        pangkat $b^N$, jika dan hanya jika setiap faktor prima $q$
        \omterm{def:b1:arith:divides}{membagi} $b$. Periksa: $\frac{1}{20}$ berakhir dalam basis
        $10$, tetapi tidak dalam basis $3$.
  \item Definisikan \emph{pemenggalannya} $s_n = \lfloor b^n x \rfloor
        / b^n$. Untuk $x = \sqrt 2$ dan $b = 10$, hitung $s_0,
        \dots, s_4$ dengan memeriksa pada setiap langkahnya bahwa dua
        kuadrat berurutan mengapit $2$ (misalnya $1.4142^2
        = 1.99996164 < 2 < 2.00024449 = 1.4143^2$), lalu periksa
        $s_n \leq \sqrt 2 < s_n + 10^{-n}$ setiap kalinya.
\end{enumerate}

\medskip
\noindent\textbf{Bagian II --- Keberadaannya, dari \omterm{def:b1:reals:bounds}{supremum}.} Tetapkan
$x \in \intco{0}{1}$ lalu tulis $A_n = \lfloor b^n x \rfloor$ dan
$d_n = A_n - b\,A_{n-1}$ untuk $n \geq 1$.
\begin{enumerate}[resume]
  \item Tunjukkan $A_0 = 0$ dan $b\,A_{n-1} \leq A_n \leq b\,A_{n-1} +
        b - 1$; lalu simpulkan bahwa setiap $d_n$ merupakan angka.
  \item Tunjukkan bahwa $s_n := A_n b^{-n}$ memenuhi
        \[
        s_n = \sum_{k=1}^{n} d_k\,b^{-k}
        \qquad\text{dan}\qquad
        s_n \leq x < s_n + b^{-n} .
        \]
  \item Buktikan $b^n \geq n + 1$ dengan induksi, lalu tunjukkan bahwa
        $(s_n)$ tidak turun dan bahwa $x = \sup_n s_n$
        \emph{(pakai \cref{prop:b1:reals:epsilon} dan
        \cref{thm:b1:reals:archimedes})}.
  \item Tunjukkan bahwa untaian $(d_n)$ bersifat sejati: karena jika $d_k = b - 1$
        untuk setiap $k > N$, hitunglah $s_n$ untuk $n > N$ lewat jumlah
        geometri yang hingga lalu bantahlah pertanyaan 6.
  \item Sebaliknya, misalkan $(e_n)_{n \geq 1}$ sebarang untaian angka
        yang sejati dan $t_n = \sum_{k=1}^n e_k b^{-k}$. Tunjukkan bahwa $y
        = \sup_n t_n$ ada, terletak di $\intco{0}{1}$, dan
        memenuhi $t_n \leq y < t_n + b^{-n}$ untuk setiap $n$
        \emph{(untuk ketaksamaan tegasnya, pakai sebuah angka $e_m \leq b
        - 2$ dengan $m > n$)}. Lalu simpulkan $\lfloor b^n y \rfloor = b^n
        t_n$, dan kemudian bahwa angka $y$, dalam pengertian
        pertanyaan 5, tepat sama dengan $e_n$.
\end{enumerate}

\medskip
\noindent\textbf{Bagian III --- Ketunggalan, urutan, geseran.}
\begin{enumerate}[resume]
  \item Rakit pertanyaan 5--9 menjadi \emph{teorema ekspansi
        $b$-adik}: bahwa \omterm{def:b1:logic:map}{pemetaan} $x \mapsto (d_n)$ dan $(e_n) \mapsto
        \sup_n t_n$ saling invers dan merupakan bijeksi antara
        $\intco{0}{1}$ dan \omterm{def:b1:logic:sets}{himpunan} untaian angka yang sejati. Khususnya
        tak ada dua untaian sejati yang berbeda yang bernilai sama.
  \item Sekarang bolehkan untaian yang tak sejati. Tunjukkan bahwa untaian dengan $e_n
        = b - 1$ untuk setiap $n > M$ (dengan $M \geq 0$ yang minimal) bernilai
        $t_M + b^{-M}$; lalu simpulkan bahwa $0.999\dots = 1$ dalam
        basis $10$, dan bahwa bilangan real yang mempunyai dua penyajian
        angka tepat berupa pecahan $b$-adik
        $m/b^N \in \intoo{0}{1}$ --- sedangkan setiap bilangan real lainnya hanya punya
        satu, bahkan di antara untaian yang tak sejati.
  \item Buktikan bahwa bijeksi pertanyaan 10 bersifat
        mengawetkan urutan bagi urutan leksikografis: yakni jika
        untaian sejati $x$ dan $y$ pertama kali berbeda pada indeks $m$,
        maka $x < y$ jika dan hanya jika $d_m < e_m$.
  \item (Lema geseran) Misalkan $x \in \intco{0}{1}$ berangka
        $(d_n)$. Tunjukkan bahwa bagian pecahan $bx$ berangka
        $(d_{n+1})_{n \geq 1}$ \emph{(hitunglah $\lfloor b^n(bx -
        A_1)\rfloor$ memakai $\lfloor u - K \rfloor = \lfloor u
        \rfloor - K$ untuk $K$ bulat)}, lalu simpulkan dengan induksi
        bahwa bagian pecahan $b^m x$ berangka
        $(d_{n+m})_{n \geq 1}$.
\end{enumerate}

\medskip
\noindent\textbf{Bagian IV --- Kerasionalan adalah keperiodikan.} Misalkan $x
= \frac pq \in \intco{0}{1}$ dalam bentuk paling sederhana dan $r_n = b^n p
\bmod q$ sisa pembagian Euclid $b^n p$ oleh
$q$.
\begin{enumerate}[resume]
  \item Tunjukkan $A_n = \dfrac{b^n p - r_n}{q}$ dan $r_n =
        (b\,r_{n-1}) \bmod q$.
  \item Tunjukkan $d_n = \Bigl\lfloor \dfrac{b\,r_{n-1}}{q}
        \Bigr\rfloor$: jadi setiap angkanya merupakan fungsi dari sisa
        sebelumnya saja. Dan inilah persis pembagian panjang pada
        Bagian I.
  \item Terapkan prinsip sarang merpati
        (\cref{cor:b1:counting:pigeonhole}) pada $r_0, \dots, r_q$
        lalu simpulkan: ekspansi setiap bilangan rasional bersifat
        periodik pada akhirnya, dengan praperiode dan periode paling besar
        $q$.
  \item Sebaliknya, andaikan angka $y \in \intco{0}{1}$
        bersifat periodik \emph{murni}: $d_{n+T} = d_n$ untuk setiap $n
        \geq 1$. Dengan memakai lema geseran dan ketunggalan pada
        pertanyaan 10, tunjukkan bahwa bagian pecahan $b^T y$
        sama dengan $y$, lalu simpulkan $(b^T - 1)\,y \in \N$: jadi $y$
        rasional dengan penyebut yang \omterm{def:b1:arith:divides}{membagi} $b^T - 1$. Periksa
        mekanismenya pada $0.(142857)$: $142857 \times 7 = 999999$.
  \item Tangani kasus periodik pada akhirnya lewat penggeseran, lalu nyatakan
        \emph{kriteria keperiodikannya}: bahwa $x \in \intco{0}{1}$
        rasional jika dan hanya jika ekspansi $b$-adik sejatinya
        periodik pada akhirnya --- dalam satu basis jika dan hanya jika dalam
        semua basis.
  \item Untuk $x = \frac 1q$ dengan $\gcd(q, b) = 1$, tunjukkan bahwa
        ekspansinya periodik murni dan bahwa periode terkecilnya adalah
        $T \geq 1$ terkecil dengan $b^T \equiv 1 \pmod q$ (yaitu
        orde perkalian $b$ modulo $q$). Periksa bahwa untuk
        $q = 7$, $b = 10$ pangkat $10$ modulo $7$ berjalan
        melewati $3, 2, 6, 4, 5, 1$: jadi ordenya $6$, yang cocok dengan pertanyaan
        1.
\end{enumerate}

\medskip
\noindent\textbf{Bagian V --- Dividen dan diagonalnya.}
\begin{enumerate}[resume]
  \item Misalkan $x^*$ bilangan real di $\intco{0}{1}$ yang angka
        basis $10$-nya bernilai $1$ pada posisi segitiga $\frac{j(j +
        1)}{2}$ ($j \geq 1$) dan $0$ di tempat lain: $x^* =
        0.101001000100001\dots$ Tunjukkan bahwa untaian angkanya
        sejati tetapi tidak periodik pada akhirnya \emph{(karena periode $T$
        akan memaksa angka satu berjarak paling jauh $T$, padahal jaraknya tumbuh)},
        lalu simpulkan bahwa $x^*$ irasional: yakni bilangan yang terbukti
        irasional lewat iramanya belaka.
  \item Tunjukkan bahwa untuk setiap basis $b$ \omterm{def:b1:logic:sets}{himpunan} $\{m/b^n : m \in \Z,
        n \in \N\}$ padat di $\R$ (yang menyamaratakan
        \cref{exo:b1:reals:8}), dan bahwa setiap bilangan rasional $\frac pq
        \in \intoo{0}{1}$ mempunyai ekspansi yang \emph{berakhir} dalam
        basis $q$. Moralnya: berakhir adalah sifat pasangan
        (bilangan, basis); sedangkan keperiodikan --- yakni kerasionalan --- bersifat
        hakiki.
  \item (Diagonal Cantor) Misalkan $k \mapsto x_k$ sebarang \omterm{def:b1:logic:map}{pemetaan} dari
        $\N^*$ ke $\intco{0}{1}$. Definisikan untaian angka $e_k =
        1$ bila angka ke-$k$ dari $x_k$ berbeda dari $1$, dan
        $e_k = 2$ bila tidak. Tunjukkan bahwa $(e_k)$ sejati, bahwa
        nilainya $y$ terletak di $\intco{0}{1}$, dan bahwa $y \neq
        x_k$ untuk setiap $k$. Lalu simpulkan: tak ada \omterm{def:b1:logic:map}{pemetaan} $\N^* \to
        \intco{0}{1}$ yang \omterm{def:b1:logic:inj}{surjektif}. (Adapun kosakata
        keterbilangan, beserta rumah yang sepantasnya bagi teorema ini, ada pada
        \cref{ch:b1:topology}.)
  \item Tunjukkan bahwa jika ekspansi sejati $x$ dan $y$ bersesuaian
        sampai indeks $n$ maka $\abs{x - y} < b^{-n}$, lalu bantahlah
        konversnya dengan $x = 0.1$, $y = 0.0999$ dalam basis $10$:
        jadi kedekatan bilangan tak memaksa kesesuaian angkanya.
        Bilangan real mana yang harus disalahkan?
  \item Jalankan Bagian IV pada $x = \frac{1}{10}$ dalam basis $b = 2$:
        hitung sisa dan angkanya sampai berdaur, lalu
        simpulkan $\frac{1}{10} = (0.0\overline{0011})_2$, dengan
        praperiode $1$ dan periode $4$. Jelaskan lewat pertanyaan 3 mengapa
        tak akan pernah ada untaian biner hingga yang sama dengan $\frac{1}{10}$
        --- yaitu alasan mengapa $0.1 + 0.2$ pada titik mengambang sebuah komputer
        tidak persis $0.3$.
  \item Sintesis. Dalam satu kalimat untuk masing-masing: di mana buktinya memakai
        (i) kelengkapan, (ii) \omterm{thm:b1:reals:archimedes}{sifat Archimedes}, (iii)
        klausa ketunggalan pada \omterm{thm:b1:reals:floor}{lantainya}, (iv) prinsip sarang
        merpati? Lalu moralnya: bahwa $\intco{0}{1}$ tersandikan dengan setia
        oleh untaian angka yang sejati, dan kerasionalan terbaca sebagai
        keperiodikan --- namun analisis lebih menyukai \omterm{def:b1:reals:bounds}{supremum} daripada
        angkanya. Mengapa? (Pikirkan tentang menjumlahkan dua untaian angka.)

\end{enumerate}
\end{problem}
